\documentclass[11pt]{article}
\usepackage{amsmath, amssymb}
\usepackage{geometry}
\usepackage{array}
\usepackage{booktabs}
\geometry{margin=0.7in}
\setlength{\parindent}{0pt}
\setlength{\parskip}{0.35em}
\renewcommand{\arraystretch}{1.15}

\newcommand{\cheatsheettitle}[2]{%
\begin{center}
{\LARGE #1}\\[0.4em]
{\large #2}
\end{center}
}


\begin{document}
\cheatsheettitle{Algebra I Functions and Modeling Cheatsheet}{Module 10}

\section*{Function Rule}
A function assigns each input exactly one output.
\[
x\rightarrow f(x)
\]
Vertical line test:
\[
\text{a graph is a function if every vertical line hits it at most once}
\]

\section*{Function Notation}
\[
f(x)=2x-7
\]
\[
f(5)=2(5)-7=3
\]
The input is the value substituted for $x$.

\section*{Domain and Range}
\[
\text{domain}=\text{set of input values}
\]
\[
\text{range}=\text{set of output values}
\]
For tables, list the $x$-values for the domain and the $y$-values for the range.

\section*{Representations}
Functions can be shown with:
\[
\text{equations},\qquad \text{tables},\qquad \text{graphs},\qquad \text{verbal descriptions}
\]
Connections:
\[
\text{table differences}\leftrightarrow \text{rate of change}
\]
\[
\text{graph intercepts}\leftrightarrow \text{starting values or zeros}
\]

\section*{Arithmetic Sequences}
\[
a_n=a_1+(n-1)d
\]
Constant difference:
\[
d=a_n-a_{n-1}
\]
Arithmetic sequences are linear functions of the term number.

\section*{Simple Exponential Growth}
\[
y=a(1+r)^t
\]
Growth:
\[
r>0
\]
Decay:
\[
y=a(1-r)^t,\qquad 0<r<1
\]
Exponential tables have a constant ratio.

\section*{Model Choice}
\begin{center}
\begin{tabular}{lll}
\toprule
Pattern & Table & Model \\
\midrule
constant difference & add same amount & linear \\
constant ratio & multiply by same amount & exponential \\
symmetry with turning point & second differences constant & quadratic \\
\bottomrule
\end{tabular}
\end{center}

\section*{Quick Checks}
\begin{itemize}
  \item Inputs cannot repeat with different outputs in a function.
  \item Linear models add by a constant amount.
  \item Exponential models multiply by a constant amount.
  \item Interpret variables with units when modeling.
\end{itemize}

\end{document}
